\notechapter{N-20251108}{Notes on MML: Vectors, the Minus-One Trick, Linear Maps, Basis Change, Kernel, and Image}

While reading the linear-algebra chapters of MML, the same objects kept moving between geometric, computational, and structural descriptions.  The first example is the word ``vector'' itself.

\section{Vectors beyond arrows}

Vectors are not only geometric arrows.  Polynomials form a vector space because they can be added and scaled.  Audio signals can likewise be represented by lists of numbers whose addition and scalar multiplication remain valid signals.

Machine-learning embeddings provide another example.  Word and sentence vectors are not literal geometric arrows, but they live in vector spaces and obey vector operations; their meaning comes from the learned representation.

The lesson is that a vector is an element of a vector space.  Geometry is only one model.

\section{The minus-one trick}

The MML notes use a practical trick for reading solutions of a homogeneous system
\[
  Ax=0,\qquad A\in\mathbb R^{k\times n}.
\]
Assume $A$ is in reduced row echelon form and has no zero rows.  Pivot columns correspond to leading ones; free columns correspond to variables that can be chosen freely.

The minus-one trick augments the row-reduced matrix into an $n\times n$ matrix by inserting rows with a single $-1$ in the diagonal positions corresponding to free variables.  Then the columns containing $-1$ on the diagonal directly give basis vectors for the null space.

The point is not magic.  It is a compact way of doing the usual free-variable substitution.  Setting one free variable to $1$ and the others to $0$ gives one kernel basis vector; the inserted $-1$ row records that choice.

\section{Example of the trick}

Consider
\[
  A=
  \begin{bmatrix}
  1&3&0&0&3\\
  0&0&1&0&9\\
  0&0&0&1&-4
  \end{bmatrix}.
\]
The free variables are in columns $2$ and $5$.  Augmenting gives
\[
  \widetilde A=
  \begin{bmatrix}
  1&3&0&0&3\\
  0&-1&0&0&0\\
  0&0&1&0&9\\
  0&0&0&1&-4\\
  0&0&0&0&-1
  \end{bmatrix}.
\]
The columns with diagonal $-1$ give
\[
  x=\lambda_1
  \begin{bmatrix}3\\-1\\0\\0\\0\end{bmatrix}
  +\lambda_2
  \begin{bmatrix}3\\0\\9\\-4\\-1\end{bmatrix}.
\]
This is the same solution one obtains by solving by free variables, but the trick makes the basis easy to read.

\section{Linear mappings}

A mapping $\Phi:V\to W$ between vector spaces is linear if
\[
  \Phi(\lambda x+\psi y)=\lambda \Phi(x)+\psi \Phi(y)
\]
for all $x,y\in V$ and scalars $\lambda,\psi$.

Important special cases are:
\begin{itemize}
  \item an isomorphism is linear and bijective;
  \item an endomorphism is a linear map $V\to V$;
  \item an automorphism is a bijective endomorphism;
  \item the identity map $\operatorname{id}_V:V\to V$ is the identity automorphism.
\end{itemize}

These notions may be summarized as “same structure,” “same domain and codomain,” and “self-isomorphism.”

\section{Complex numbers as a real vector space}

The example
\[
  \Phi:\mathbb R^2\to\mathbb C,\qquad
  \Phi(x_1,x_2)=x_1+ix_2
\]
is linear over $\mathbb R$:
\[
  \Phi(x+y)=\Phi(x)+\Phi(y),\qquad
  \Phi(\lambda x)=\lambda \Phi(x).
\]
It is also bijective, so it is an isomorphism of real vector spaces.

The note marks this as a homomorphism.  The important distinction is that it preserves the vector-space addition and scalar multiplication.  It is not claiming that every additional structure is automatically preserved.

\section{Finite-dimensional isomorphism theorem}

The theorem quoted from Axler says that finite-dimensional vector spaces $V$ and $W$ are isomorphic iff
\[
  \dim V=\dim W.
\]
This is one of the most important simplifications in finite-dimensional linear algebra: after choosing bases, all $n$-dimensional vector spaces over the same field look like $\mathbb F^n$.

But this isomorphism is not canonical unless a basis is chosen.  The abstract vector space and its coordinate representation are related but not identical.

\section{Basis change}

Let $\Phi:V\to W$ be linear.  Suppose $V$ has old basis
\[
  B=(b_1,\ldots,b_n)
\]
and new basis
\[
  \widetilde B=(\widetilde b_1,\ldots,\widetilde b_n).
\]
Similarly, let $W$ have old basis
\[
  C=(c_1,\ldots,c_m)
\]
and new basis
\[
  \widetilde C=(\widetilde c_1,\ldots,\widetilde c_m).
\]
If $A_\Phi$ is the matrix of $\Phi$ with respect to $B,C$, and $\widetilde A_\Phi$ is the matrix with respect to $\widetilde B,\widetilde C$, then
\[
  \widetilde A_\Phi=T^{-1}A_\Phi S.
\]
Here $S$ converts coordinates from the new basis of $V$ to the old basis of $V$, and $T$ converts coordinates from the new basis of $W$ to the old basis of $W$.

The coordinate-change diagram expresses exactly this procedure: translate the input coordinates into the original basis, apply $A_\Phi$, and then translate the output coordinates into the new output basis.

\section{Equivalence and similarity}

Two rectangular matrices $A,\widetilde A\in\mathbb R^{m\times n}$ are equivalent if there exist invertible matrices $S\in\mathbb R^{n\times n}$ and $T\in\mathbb R^{m\times m}$ such that
\[
  \widetilde A=T^{-1}AS.
\]
This corresponds to changing the basis in both domain and codomain.

Two square matrices $A,\widetilde A\in\mathbb R^{n\times n}$ are similar if
\[
  \widetilde A=S^{-1}AS.
\]
This corresponds to changing basis for a linear map $V\to V$ using the same coordinate change in domain and codomain.

The distinction is:
\begin{center}
\begin{tabular}{lll}
\toprule
Concept & Situation & Invariant meaning \\
\midrule
Equivalence & linear map $V\to W$ & change bases in domain and codomain \\
Similarity & linear map $V\to V$ & same operator in a new basis \\
\bottomrule
\end{tabular}
\end{center}

\section{Kernel and image}

For a linear map $\Phi:V\to W$, the kernel is
\[
  \ker \Phi=\{v\in V:\Phi(v)=0_W\},
\]
and the image is
\[
  \operatorname{Im}\Phi=\{w\in W:\exists v\in V,\ \Phi(v)=w\}.
\]
The kernel records what is collapsed to zero.  The image records what outputs are reachable.

The source diagram shows the domain on the left, codomain on the right, a yellow region for the kernel, and a blue/yellow region for the image.  The note also records the injectivity criterion:
\[
  \Phi\text{ is injective}\quad\Longleftrightarrow\quad\ker \Phi=\{0_V\}.
\]

\section{Categorical side remarks}

The margin notes compare kernel, image, and cokernel with categorical language.  In ordinary linear algebra,
\[
  \ker f\to A\to B
\]
means the kernel is a subobject of the domain.  The image is a subobject of the codomain.  The cokernel, when defined, is a quotient-like object on the codomain side.

This is a useful preview: linear algebra already contains many category-theoretic patterns, even before they are named abstractly.

\section{The general pattern}

This note links computational linear algebra with structural thinking.  Row reduction gives the null space; the minus-one trick packages free-variable substitution.  Linear maps become matrices only after choosing bases.  Basis change explains why the same map can have different matrices.  Kernel and image describe what a map destroys and what it reaches.  These ideas are also the language behind machine-learning embeddings, where vectors need not be arrows but still live in vector spaces with meaningful linear structure.

\section{Vectors as elements of a linear structure}
A vector is not essentially an arrow; it is an element of a vector space.  Coordinates appear only after choosing a basis.  The same vector has different coordinate columns in different bases.
\section{The minus-one trick as a kernel graph}

The phrase ``minus-one trick'' should be understood carefully.  In this note it is not primarily a homogeneous-coordinate trick and it is not an affine shift.  It is a compact way of reading the kernel of a row-reduced homogeneous system.

\begin{proposition}[Kernel graph form of an RREF system]
After possibly reordering variables, suppose a homogeneous system has row-reduced matrix
\[
  \begin{bmatrix} I_r&A\end{bmatrix}.
\]
Write the unknown vector as
\[
  x=(x_p,x_f),
\]
where \(x_p\in k^r\) are pivot variables and \(x_f\in k^{n-r}\) are free variables.  Then the solution space is
\[
  \ker\begin{bmatrix} I_r&A\end{bmatrix}
  =
  \left\{
  \begin{pmatrix}-Au\\ u\end{pmatrix}
  :u\in k^{n-r}
  \right\}.
\]
Thus the kernel is the graph of the linear map
\[
  u\longmapsto -Au.
\]
\end{proposition}

\begin{proof}
The equation is
\[
  x_p+A x_f=0.
\]
Hence \(x_p=-A x_f\).  If we write \(u=x_f\), every solution has the displayed form, and every displayed vector clearly satisfies the equation.
\end{proof}

The inserted \(-1\) rows in the practical trick are a bookkeeping device for the free variables.  Choosing \(u=e_j\), the \(j\)-th standard basis vector of the free-variable space, gives one kernel basis vector:
\[
  \begin{pmatrix}-Ae_j\\ e_j\end{pmatrix}.
\]
So the trick is best read as a way of writing the graph basis of a kernel.

\begin{remark}[Do not confuse two uses of extra coordinates]
Homogeneous coordinates do turn affine transformations into linear transformations one dimension higher.  That is a real and useful idea.  But the minus-one trick in this row-reduction context is different: it does not linearize an affine map; it packages the free-variable choices in a homogeneous linear system.
\end{remark}
\section{Kernel-image factorization}
For a linear map \(T:V\to W\), the kernel records lost directions and the image records reachable directions.  Rank-nullity says
\[
\dim V=\dim\ker T+\dim\operatorname{im}T.
\]
\section{Coordinates, kernels, and affine linearization}
The review should distinguish vectors from coordinates, affine maps from linear maps, and a matrix from the linear transformation it represents.  Kernel, image, and basis change are the structural core.

\section{Exact sequences behind rank-nullity}

For a linear map $T:V\to W$, the kernel and image fit into the exact sequence
\[
  0\to \ker T \to V \xrightarrow{T} \operatorname{im}T \to 0.
\]
Exactness means that the image of each arrow equals the kernel of the next.  Taking dimensions gives
\[
  \dim V=\dim \ker T+\dim \operatorname{im}T.
\]
This is rank-nullity.

The elementary theorem is therefore the dimension shadow of an exact sequence.  Exact sequences become one of the central languages of algebra, topology, and geometry.

\section{Quotient description of the image}

The first isomorphism theorem says
\[
  V/\ker T\cong \operatorname{im}T.
\]
The quotient collapses all vectors that differ by a null vector.  What remains is precisely the information visible after applying $T$.

This gives a clean meaning to free variables and solution spaces.  The kernel is what the map forgets; the image is what survives.

\section{Cokernel and obstructions}

The cokernel of $T:V\to W$ is
\[
  \operatorname{coker}T=W/\operatorname{im}T.
\]
It measures the failure of $T$ to be onto.  A system
\[
  Tx=b
\]
is solvable exactly when $b$ maps to zero in the cokernel.

Thus kernel measures non-uniqueness, while cokernel measures obstruction to existence.  This is the conceptual form of the elementary distinction between infinitely many solutions and no solution.

\section{Kernel, image, and quotient survive basis changes}

The minus-one trick, basis conversion, and kernel/image computations in the note are all examples of one structure: a linear map forgets some directions, preserves some information, and may fail to reach some outputs.  Exact sequences and quotients express this without depending on a chosen matrix.

\section{Row reduction as choosing a splitting}

When a matrix is row-reduced, the pivot variables are expressed in terms of free variables.  Abstractly, this chooses a splitting of the domain into a complement of the kernel plus the kernel:
\[
  V\cong U\oplus \ker T.
\]
The map $T$ is injective on the chosen complement $U$ and maps $U$ isomorphically onto $\operatorname{im}T$.

This splitting is not canonical: different row-reduction choices or different bases may choose different complements.  The exact sequence
\[
  0\to\ker T\to V\to\operatorname{im}T\to0
\]
is canonical; a particular parametric solution is a coordinate choice.

\section{Cokernel as the missing equation space}

For $T:V\to W$, the cokernel $W/\operatorname{im}T$ measures which right-hand sides are not reachable.  In matrix terms, left null vectors produce compatibility conditions.  If $y\in\ker T^T$, then any solvable equation $Tx=b$ must satisfy
\[
  y^Tb=0.
\]
Thus the left null space is the dual of the obstruction space.

This gives a precise interpretation of consistency conditions in linear systems: they are not extra tricks, but equations defining whether $b$ vanishes in the cokernel.
